\documentclass[11pt,a4paper]{article}

\usepackage[margin=2.5cm]{geometry}
\usepackage{setspace}
\usepackage{graphicx}
\usepackage{float}
\usepackage[most]{tcolorbox}
\usepackage{graphicx}
\usepackage{tabularx}
\usepackage{array}
\usepackage{tikz}
\usepackage{booktabs}
\usetikzlibrary{arrows.meta, positioning, shapes.geometric}
\setstretch{1.15}

\title{Progetto Signal, Image and Video}
\author{}
\date{}

\begin{document}
\maketitle

\begin{enumerate}
    \item Dataset iniziale
    \begin{itemize}
        \item Shape, classi e distribuzione (quanti sample per classe, quanti soggetti)
        \item Verifica valori mancanti sui landmark grezzi
    \end{itemize}
    \item Landmark selezionati
    \begin{itemize}
        \item Mostra quali landmark usi e perché
        \item Distribuzione della visibility per landmark critico (spalle, anche) — per giustificare le soglie che usi dopo
    \end{itemize}
    \item Midpoint construction
    \begin{itemize}
        \item Verifica disponibilità dei midpoint
        \item Distribuzione delle sorgenti di head\_center
    \end{itemize}
    \item Feature engineering
    \begin{itemize}
        \item Descrizione breve di cosa misura ogni feature
        \item Distribuzione di ogni feature per classe
        \item Valori NaN per feature
        \item Normalizzazione \textit{(eps, wrap\_angle\_90, angle\_diff\_deg)}
    \end{itemize} 
    \item Dataset finale
    \begin{itemize}
        \item Shape finale, 8 primary + 1 support
        \item Distribuzione classi nel dataset pulito
        \item Split soggetti train/test — no data leakage
    \end{itemize}
\end{enumerate}

\newpage
\section{Dataset Analysis}
\subsection{Dataset iniziale}
Il dataset grezzo è stato ottenuto acquisendo sequenze video di \textbf{3 soggetti}
che eseguono variazioni posturali del tronco superiore. L'estrazione dei landmark
3D è stata effettuata tramite \textit{MediaPipe Pose}, che restituisce per ogni
frame le coordinate $(x, y, z)$ di 33 keypoint corporei.
\end{document}
